\section{Relation to Prior Work}\label{sec:related-work}

\subsection{Free completion and partial algebras}
The free compatible completion underlying \cref{sec:free-completion} belongs to
an established line of work on partial algebras and Schmidt-type completions
\cite{ClimentVidalCosmeLlopez2023}.  Diaconescu's many-sorted encoding provides
a closely related partial-to-total perspective \cite{Diaconescu2009}.  These
results motivate the substrate but are not novelty claims of the present paper.

\subsection{Recognizable trees and finite separators}
For finite bases and finite signatures, the normalized ground-equational
presentation lies within classical tree-automata and recognizability theory
\cite{Kozen1977,Brainerd1969,DauchetTison1990,ComonEtAl2008}.  Finite selected
subterm closures plus a sink state are familiar constructions.  The relative
feature studied here is the requirement that every observer preserve an
arbitrary partial base pointwise while only the complement over that base is
budgeted.  This permits infinite codomains with finite external complexity.

There is a concrete obstruction to reducing the simultaneous finite-pattern
statement to the usual ``separate pairs and take their product'' argument.
The obstruction is already visible at carrier level.  For an infinite base
$A$, each one-point extension $A\hookrightarrow A\sqcup\{\star\}$ has finite
relative complement.  But the Cartesian product of two such extensions, with
$A$ embedded diagonally, contains all mixed points $(a,\star)$ outside the old
image.  Hence its complement is infinite.  This failure of product closure is
Lean-checked for $A=\mathbb N$ in
\texttt{ResolutionRelativeProductObstruction.lean}.  Thus abstract
finite-family discrimination may be classical, but the standard product
construction does not establish it inside the restricted observer class used
here when the preserved base is infinite; the direct finite-pattern
construction is doing additional work.

\subsection{Filtered and profinite completions}
Cauchy completion of separated filtrations and factorial synchronization of
finite deterministic dynamics are standard.  Profinite completions organize
observations by finite quotients or finite recognizing algebras
\cite{DaveyGouveiaHaviarPriestley2011,ChenAdamekMiliusUrbat2016,Moreau2024}.
Resolution stages differ because the preserved base may be infinite.  The
compression theorem deliberately isolates the standard dynamic ingredient:
only the trajectory relevant to the chosen unary family needs a finite code.
The new structural interaction is with the relative finite-pattern escape
observer that prevents convergence to any finite generated candidate.

\subsection{Finite complements and Rees-index analogies}
Finite-complement extensions and residual finiteness are classical in semigroup
theory \cite{CainMaltcev2013,RuskucThomas1998}.  The analogy is useful but not
literal here.  In a genuinely partial Resolution algebra, the embedded old
carrier is not closed under the generated total operations; by
\cref{prop:old-image-closed}, closure occurs exactly when the evaluator was
already total.  Hence the observer class is better described as relative to a
pointwise-fixed partial base than as an ordinary finite-Rees-index extension.

\subsection{Tree metrics}
Classical infinite-tree completions are based on prefix depth
\cite{ArnoldNivat1980,Courcelle1983}.  Resolution uses semantic finite-tag
observation instead.  The formal comb counterexample in \cref{sec:completion}
shows that the two Cauchy notions differ: a sequence may be depth-Cauchy while a
fixed low-budget Resolution observer continues to distinguish adjacent terms.

\subsection{Contribution and scope}
The contribution claimed by this paper is concentrated in the following
package:
\begin{enumerate}
  \item an intrinsic finite-complement observer class over a pointwise-fixed,
        possibly infinite partial base, together with finite-pattern
        realization;
  \item the induced observational filtration and its trajectory-level
        compression theorem;
  \item the Resolution-specific finite-pattern escape mechanism and the
        resulting compression--escape properness theorem;
  \item the exact finite-base criterion, infinite-base old-fixing instance,
        sharper unbounded-rank witnesses, and exact equational conservativity.
\end{enumerate}
The factorial-periodicity argument itself is treated as standard machinery.
Likewise, the many-sorted Diaconescu development is a verified compatibility and
provenance layer rather than an independent novelty claim.
